\section{Four frequencies in a wide arithmetic progression}
\label{sec:four}

For positive integers $N,m$, let $M_m^{(4)}(N)$ be the supremum of lengths
of intervals on which a nonzero real trigonometric polynomial with spectrum
contained in
\[
 \{\pm N,\ \pm(m+1)N,\ \pm(2m+1)N,\ \pm(3m+1)N\}
\]
does not vanish. Length is measured in $\mathbb R/\mathbb Z$ units. The result
below extends the integer-ratio three-frequency family in the repository's
incoming \texttt{three\_frequency\_gaps.tex}. It is an infinite family of
four-frequency cases of Kozma--Oravecz Question 1~\cite{KO2003}, not a solution for every
arithmetic progression or every number of frequencies. It is unconditional.

For $0\le e\le1/3$ put
\[
 h_e(s)=\begin{cases}\sin(es)/e,&e>0,\\s,&e=0,\end{cases}
 \qquad L_e=\begin{cases}\pi/(2e),&e>0,\\+\infty,&e=0.\end{cases}
\]
The entire extension at $e=0$ is intended in all derivatives. Define
\begin{equation}\label{eq:fourF}
 F_e(s)=h_e(s)+A(e)\cos((1+e)s)+B(e)\sin((2+e)s)
                  +C(e)\cos((3+e)s).
\end{equation}
The coefficients and the four contact parameters $x<0<u<v<w<L_e$ are
specified by the seven equations
\begin{equation}\label{eq:fourcontacts}
 F_e(u)=F_e(v)=F_e(w)=F_e(x)=0,
 \qquad F'_e(u)=F'_e(v)=F'_e(w)=0,
\end{equation}
together with the conditions in the certificate lemma below.

\begin{theorem}[Four-frequency gap formula]\label{thm:four}
For every positive integer $N$ and integer $m\ge2$,
\[
 N M_m^{(4)}(N)=
 \begin{cases}
  1/2,&m\text{ even},\\[2pt]
  1/2-x(1/m)/(\pi m),&m\text{ odd}.
 \end{cases}
\]
For odd $m$, the supremum is not attained even by an open zero-free interval.
For these odd $m$, on the centered critical closed interval, the cone of nonzero nonnegative
polynomials with the stated spectrum is one-dimensional. There is precisely
one annihilating probability measure on that interval: it has eight atoms,
each of strictly positive mass. Here support length means the length of a shortest closed arc containing
the support. For odd $m$, every annihilating measure of minimum support length is a
translate of this one.

There is a real analytic branch of $(A,B,C,u,v,w,x)$ on a neighborhood of
$[0,1/3]$. Consequently the odd-$m$ gap has a convergent expansion in $1/m$
near zero. Its first terms are
\begin{align*}
 N M_m^{(4)}(N)&=\frac12+\frac{c_1}{m}+\frac{c_2}{m^2}+O(m^{-3}),\\
 c_1&=-\frac{x(0)}\pi=0.45356051377648\ldots,\\
 c_2&=-\frac{x'(0)}\pi=-0.18321967877208\ldots.
\end{align*}
In particular,
\[
 0.63147839169477181<M_3^{(4)}(1)<0.63147839169477182.
\]
\end{theorem}

\begin{lemma}[Continuum certificate]\label{lem:fourcertificate}
For every $e\in[0,1/3]$ there is a solution of
\eqref{eq:fourcontacts}, with $-2<x<0<u<v<w<5$, such that:
\begin{enumerate}
\item $w<L_e$, and $F_e\ge0$ on $[x,L_e]$, with zeros exactly
$x,u,v,w$. Moreover $F'_e(x)>0$ and
$F''_e(u),F''_e(v),F''_e(w)>0$.
\item Let the columns of the $3\times4$ matrix $V$ be
\[
 (\cos((1+e)s),\ \sin((2+e)s),\ \cos((3+e)s))^\top,
 \qquad s=x,u,v,w.
\]
All four signed maximal minors $(-1)^j\det V_{\widehat j}$,
$j=0,1,2,3$, are negative.
\item The $4\times4$ matrix with columns
\[
 (\cos(es),\ \sin((1+e)s),\ \cos((2+e)s),\ \sin((3+e)s))^\top,
 \qquad s=x,u,v,w,
\]
has positive determinant.
\end{enumerate}
The Jacobian of the seven equations with respect to
$(A,B,C,u,v,w,x)$ is nonsingular at every solution provided by the certificate.
\end{lemma}

\begin{proof}[How the finite certificate proves the lemma]
In \texttt{checks/spectral}, run:
\begin{verbatim}
python3 -S four_frequency_certificate.py replay
\end{verbatim}
The accompanying
JSON data partition $[0,1/3]$ into 1127 rational cells. Each cell includes a
rational center $z_0$, a positive rational radius $r$, and a rational
$7\times7$ preconditioner $R$. These are proposals, not trusted conclusions.
Let $g(e,z)$ be the vector in \eqref{eq:fourcontacts}, in the script's order,
and set $Z=z_0+[-r,r]^7$. The checker bounds every component of
$R g(e,z_0)$ and every entry of $I-RD_zg(e,z)$ on $E\times Z$ and verifies
\[
 q_i:=\sum_j\sup_{E\times Z}|(I-RD_zg)_{ij}|<1,
 \qquad
 \sup_E |(Rg(e,z_0))_i|+r q_i<r.
\]
For each $e\in E$ the map $z\mapsto z-Rg(e,z)$ is a strict contraction
of $Z$ into its interior. Its unique fixed point is a root of $g$:
$RD_zg$ is invertible by the norm bound, so $R$ is invertible. The same
bound proves nonsingularity of $D_zg$ at the root.

The checker verifies $F'>0$ on a neighborhood of the interval enclosure of
$x$ and $F''>0$ near each of the other three contact enclosures. The known
zero and derivative equations give exact local positivity there. Between
these neighborhoods it uses interval bisection to check $F>0$ through
$\min(9,L_e)$. In a parameter cell with lower endpoint $e_->0$, the
compact check continues to $\min(9,11/(7e_-))$; this covers the needed range
because $\pi<22/7$. The zero lower endpoint is interpreted as an infinite
cutoff. For $9\le s\le L_e$, concavity of sine gives
\[
 h_e(s)\ge 2s/\pi\ge18/\pi>63/11>|A|+|B|+|C|,
\]
where the last inequality is checked on every root box. This proves positivity
on the whole physical interval, including its unbounded limit at $e=0$.
The minor signs and the last determinant are checked directly with interval
arithmetic.

Replay uses only exact integers and rational endpoints. All elementary
operations round outwards to a $10^{-24}$ grid. Sine and cosine are evaluated
at the midpoint divided by sixteen, with Taylor degrees $49$ and $48$,
then four double-angle steps. The remainder is dominated by
$3^{50}/50!<10^{-35}$ on the checked reduced domain, and a grid unit is
added before doubling. The original argument radius is added using the
Lipschitz constant one. The function $\operatorname{sinc} y$ is evaluated
through degree $60$; $4^{62}/63!<10^{-40}$ bounds its remainder on the
checked domain, and its interval radius uses the bound
$|\operatorname{sinc}'y|\le1/2$. Thus the generation routine's floating
point root proposals are not assumptions of the proof.
As a separate check of the first nonresonant member,
\texttt{verify\_m3\_polynomial.py} works in the algebraic coordinate
$y=\cos\theta$ with
\[
 T_1(y)+aT_4(y)+bT_7(y)+cT_{10}(y).
\]
It certifies the seven contact equations, positivity on the whole interval,
positive annihilating weights and the odd-mode determinant using integer
Chebyshev polynomials. This verification uses neither trigonometric
evaluation nor the continuum witness. It therefore supplies an independent
formulation of the $\{1,4,7,10\}$ case.
\end{proof}

\begin{proof}[Uniqueness of the certified branch]
Write $d_j=(-1)^j\det V_{\widehat j}$, with columns indexed by
$(x,u,v,w)$. The four negative signed minors give the positive normalized
null vector
\[
 \lambda_j=\frac{d_j}{d_0+d_1+d_2+d_3}>0,
 \qquad \sum_{j=0}^3\lambda_j=1.
\]
Denote its entries by $\lambda_x,\lambda_u,\lambda_v,\lambda_w$.
Since $F$ vanishes at the four nodes, this vector also annihilates $h_e$.
If two certified solutions had endpoints $x_1<x_2<0$, the first polynomial
would be nonnegative on all four nodes of the second and strictly positive
at $x_2$. Its integral against the second positive rule would be positive,
contrary to annihilation. For equal endpoints the same argument forces the
four contact sets to agree. The nonzero minor obtained by removing $x$
shows that the equations at $u,v,w$ uniquely determine $A,B,C$.

The analytic implicit function theorem applies at every root, by the
nonsingular Jacobian. Solutions from adjacent parameter cells agree at
their shared endpoint by this uniqueness. Local uniqueness from the
implicit function theorem glues the analytic germs. Compactness yields an
open neighborhood of $[0,1/3]$ on which the branch is real analytic.
\end{proof}

\begin{proof}[Proof of Theorem~\ref{thm:four}]
Dilation reduces to $N=1$. When $m$ is even all four frequencies are odd,
so translating by $1/2$ changes the sign of every allowed polynomial.
No zero-free interval can have length greater than $1/2$, and the lowest
cosine has a zero-free open interval of exactly that length.

For odd $m$, set $e=1/m$ and $\theta=2\pi t$. On the positive half of
the centered critical interval let
\[
 s=m(\pi/2-\theta),\qquad
 \theta_* =\pi/2-x(e)/m.
\]
The four terms in \eqref{eq:fourF} are, up to nonzero constant factors,
the allowed cosine modes:
\[
 h_e(s)=m\cos\theta,
 \quad \cos((1+e)s)=(-1)^{(m+1)/2}\cos((m+1)\theta),
\]
\[
 \sin((2+e)s)=-\cos((2m+1)\theta),
 \quad \cos((3+e)s)=(-1)^{(3m+1)/2}\cos((3m+1)\theta).
\]
The certificate therefore gives an even allowed polynomial nonnegative on
$[-\theta_*,\theta_*]$, with precisely the eight zeros
\[
 \pm\theta_s,\qquad \theta_s=\pi/2-s/m,
 \qquad s=x,u,v,w.
\]
Place mass $\lambda_s/2$ at each of these two points. Cosine moments vanish
by the signed-minor construction and sine moments vanish by reflection.
Six atoms lie strictly inside the critical interval. Any allowed polynomial
strictly positive on its open interior has nonnegative endpoint values;
integrating against this measure is therefore strictly positive, a
contradiction. A zero-free polynomial has constant sign, so the same
argument applies after changing its sign. This excludes the critical open
gap as well as every larger gap.

For $x<a<0$, the function $F_e+\delta h_e$ is strictly positive on
$[a,L_e]$ for sufficiently small $\delta>0$. On $[a,0]$ the original
function has a strictly positive minimum, while on $(0,L_e]$ the added
term is positive; $F_e(0)>0$ as well. This is still an allowed cosine
polynomial. Its zero-free interval length tends to $\theta_*/\pi$, proving
the formula and the nonattainment assertion.

Any annihilating probability measure on the critical interval must be
supported on those eight zeros by integration of $F_e$. The paired masses
are uniquely determined by the rank-three projected matrix and mass one.
The differences within the four pairs obey the four sine equations. Up to
row signs and nonzero constants their coefficient matrix is the positive-
determinant matrix in the lemma. Thus all differences vanish and the
measure is unique. Conversely any nonnegative allowed polynomial on this
interval vanishes at every contact. Its odd part is zero by the same
invertible sine matrix, and its cosine coefficients are unique up to scale
by the rank-three condition. This proves the cone assertion. Centering
shows uniqueness up to translation for minimum support measures; smaller
support is excluded by strict positivity on every shorter arc.

Finally expand the analytic function $x(e)$ at zero. The constant values
are specified exactly by \eqref{eq:fourcontacts}; the derivative is
\[
 z'(0)=-\bigl(D_zg(0,z(0))\bigr)^{-1}\partial_e g(0,z(0)),
 \qquad z=(A,B,C,u,v,w,x).
\]
This gives the stated expansion. The decimal figures are numerical
evaluations of the exact branch definition, and are separately enclosed
by the accompanying endpoint calculations.
\end{proof}
